% оформление.tex — оформление «Клетка»: тетрадная клетка с обложки, φ, вкладки разделов.
% Только внешний вид. Логика (ссылки, длинные темы, сборка) — в shpora.sty.
% Подключается автоматически в конце преамбулы.

% ======================================================================
% 1. Цвета страницы
% ======================================================================
\colorlet{pfAkcent}{pfBrendTemn}             % коды-ссылки в тексте, «Пример»
\definecolor{pfRamka}{RGB}{196,218,246}      % рамки плашек
\definecolor{pfFon}{RGB}{236,244,255}        % полоса заголовка темы
\definecolor{pfRazborRamka}{RGB}{222,206,250}
\definecolor{pfRazborFon}{RGB}{245,239,255}
\definecolor{pfSeryj}{RGB}{110,118,130}

% Плашка-код (Ч1.3, Р11-01): \pfChip[цвет]{текст}
\renewcommand\pfChip[2][pfBrend]{\tcbox[on line, colback=#1, colupper=white, boxrule=0pt,
  arc=1mm, boxsep=0pt, left=1.2mm, right=1.2mm, top=0.75mm, bottom=0.6mm]%
  {\pfBrendZhirnyj\fontsize{8}{8}\selectfont #2}}

% ======================================================================
% 2. Клетка
% ======================================================================
\newcommand\pfKletka[4]{% x1,y1,x2,y2 в мм от левого нижнего угла страницы
  \begin{scope}[shift={(current page.south west)}]
    \clip (#1mm,#2mm) rectangle (#3mm,#4mm);
    \draw[pfKletka, line width=0.4pt, step=5mm] (0,0) grid (150mm,212mm);
  \end{scope}}

% ======================================================================
% 3. Колонтитулы
% ======================================================================
% Вкладка раздела на правом краю: место по высоте = номер раздела
\def\pf@vkladka{\ifcsdef{pf@vk@\pfRazdel}{\csuse{pf@vk@\pfRazdel}}{-1}}
\csdef{pf@vk@Числа и вычисления}{0}
\csdef{pf@vk@Алгебра}{1}
\csdef{pf@vk@Тригонометрия}{2}
\csdef{pf@vk@Функции и графики}{3}
\csdef{pf@vk@Производная}{4}
\csdef{pf@vk@Планиметрия}{5}
\csdef{pf@vk@Векторы и координаты}{6}
\csdef{pf@vk@Стереометрия}{7}
\csdef{pf@vk@Вероятность и статистика}{8}
\csdef{pf@vk@Модели: текстовые и экономические задачи}{9}
\csdef{pf@vk@Практика}{10}
\newcommand\pfVkladkaSsylka{}   % куда ведёт вкладка (начало главы / задания)
\colorlet{pfVkladkaCvet}{pfBrend} % теория — голубая, практика — фиолетовая

% Логотип и иконки собираются заранее в боксы (вложенный tikz в колонтитуле ломает позицию)
\newsavebox\pf@boxLogo \newsavebox\pf@boxSoc
\AtBeginDocument{\sbox\pf@boxLogo{\href{\pfSajt}{\pfLogoSlovo{4.6mm}}}%
  \sbox\pf@boxSoc{\pfSocseti{3.8mm}}}

\fancyhf{}
\renewcommand{\headrulewidth}{0pt}
\fancyhead[C]{\begin{tikzpicture}[remember picture, overlay]
  \pfKletka{0}{200}{150}{212}
  \begin{scope}[shift={(current page.south west)}]
    \draw[pfRamka, line width=0.5pt] (0,200mm) -- (150mm,200mm);
    \node[anchor=base west, inner sep=0pt] at (11mm,204.6mm)
      {\usebox\pf@boxLogo};
    \node[anchor=base east, inner sep=0pt, text=pfSeryj] at (139mm,204.6mm)
      {\pfBrendShrift\fontsize{7}{7}\selectfont \pfKolontitul};
    \ifnum\pf@vkladka>-1\relax
      \pgfmathsetlengthmacro\pf@vy{184mm - \pf@vkladka*14.5mm}
      \fill[pfVkladkaCvet, rounded corners=1.6mm] (143.2mm,\pf@vy-5mm) rectangle (152mm,\pf@vy+5mm);
      \node[text=white, inner sep=0pt] at (146.5mm,\pf@vy)
        {\hyperlink{\pfVkladkaSsylka}{\pfBrendZhirnyj\fontsize{6.5}{6.5}\selectfont\pfGlavaKod}};
    \fi
  \end{scope}
\end{tikzpicture}}
\fancyfoot[C]{\begin{tikzpicture}[remember picture, overlay]
  \begin{scope}[shift={(current page.south west)}]
    \draw[pfRamka, line width=0.5pt] (11mm,11mm) -- (139mm,11mm);
    \node[anchor=west, inner sep=0pt] at (11mm,6.2mm) {\usebox\pf@boxSoc};
    \node[fill=pfBrend, text=white, rounded corners=1.4mm, minimum width=9mm,
      minimum height=5mm, inner sep=0pt] at (75mm,6.2mm)
      {\pfBrendZhirnyj\fontsize{8}{8}\selectfont\thepage};
    \node[anchor=east, inner sep=0pt] at (139mm,6.2mm)
      {\hyperlink{oglavlenie}{\color{pfBrendTemn}\pfBrendShrift\fontsize{7.5}{7.5}\selectfont Оглавление\,$\uparrow$}};
  \end{scope}
\end{tikzpicture}}
\fancypagestyle{plain}{}
\fancypagestyle{pfPustaya}{\fancyhf{}}

% ======================================================================
% 4. Тема, длинная тема, разбор
% ======================================================================
\tcbset{
  pfTema/.style={enhanced, breakable=false,
    colback=white, colframe=pfRamka, boxrule=0.7pt, arc=2.2mm,
    left=3mm, right=3mm, top=2mm, bottom=2mm, halign upper=flush left,
    colbacktitle=pfFon, coltitle=black, toptitle=1.7mm, bottomtitle=1.5mm, lefttitle=2.2mm,
    titlerule=0pt, fonttitle=\bfseries\raggedright,
    title={\pfChip{#1}\enspace\csuse{pf@title@#1}},
    segmentation style={draw=pfRamka, dashed, line width=0.5pt},
    fontlower=\footnotesize\color{pfSeryj}, halign lower=flush left,
    before skip=4mm, after skip=4mm},
  % длинная тема: разрывается, на продолжении — заголовок «(продолжение)», внизу части — «дальше →»
  pfTemaDlin/.style={breakable, pad at break*=2mm,
    title after break={\pfChip{#1}\enspace\csuse{pf@title@#1}\
      {\mdseries\color{pfSeryj}(продолжение)}},
    overlay unbroken and first={},
    overlay first={\node[anchor=south east, inner sep=0pt, text=pfSeryj]
      at ([xshift=-3mm, yshift=1.6mm]frame.south east)
      {\pfBrendShrift\fontsize{7}{7}\selectfont продолжение на следующей странице $\rightarrow$};},
    overlay middle={\node[anchor=south east, inner sep=0pt, text=pfSeryj]
      at ([xshift=-3mm, yshift=1.6mm]frame.south east)
      {\pfBrendShrift\fontsize{7}{7}\selectfont продолжение на следующей странице $\rightarrow$};},
    bottomrule at break=0.7pt, toprule at break=0.7pt,
    bottomsep at break=6mm},
  pfRazbor/.style={enhanced, breakable,
    colback=white, colframe=pfRazborRamka, boxrule=0.7pt, arc=2.2mm,
    left=3mm, right=3mm, top=2mm, bottom=2mm, halign upper=flush left,
    colbacktitle=pfRazborFon, coltitle=black, toptitle=1.7mm, bottomtitle=1.5mm, lefttitle=2.2mm,
    titlerule=0pt, fonttitle=\bfseries\raggedright,
    title={\pfChip[pfPurple]{#1}\enspace\csuse{pf@title@#1}},
    title after break={\pfChip[pfPurple]{#1}\enspace\csuse{pf@title@#1}\
      {\mdseries\color{pfSeryj}(продолжение)}},
    before skip=4mm, after skip=4mm}}
\renewcommand\pfMetaOpiraetsya{{\pfMedium Опирается на:}}
\renewcommand\pfMetaGde{{\pfMedium Где встречается:}}
\renewcommand\pfMetaZadaniya{{\pfMedium Задания ЕГЭ:}~№}

% Ответ в разборе — плашкой
\renewcommand\Otvet[1]{\par\smallskip\tcbox[on line, colback=pfRazborFon, colframe=pfRazborFon,
  boxrule=0pt, arc=1mm, boxsep=0pt, left=1.5mm, right=1.5mm, top=1mm, bottom=1mm]%
  {{\color{pfPurple}\pfMedium Ответ:}\ #1}\par}

% ======================================================================
% 5. Начало главы и задания: плашка в клетку + список тем
% ======================================================================
\newcommand\pf@plashka[4]{% #1 надпись сверху, #2 крупный код, #3 название, #4 список
  \begin{tcolorbox}[enhanced, colback=white, colframe=pfRamka, boxrule=0.7pt, arc=2.2mm,
    left=4mm, right=4mm, top=3mm, bottom=3mm, before skip=0pt, after skip=3mm,
    underlay={\begin{scope}\clip[rounded corners=2.2mm] (frame.south west) rectangle (frame.north east);
      \draw[pfKletka, step=5mm, line width=0.4pt] (frame.south west) grid (frame.north east);\end{scope}}]
    {\color{pfSeryj}\pfBrendShrift\fontsize{7.5}{9}\selectfont #1}\par\vspace{2.5mm}
    {\color{pfVkladkaCvet}\pfBrendZhirnyj\fontsize{30}{30}\selectfont #2}\hspace{3mm}%
    \parbox[b]{0.66\linewidth}{\raggedright\pfBrendZhirnyj\fontsize{13}{15}\selectfont #3\vspace{0.6mm}}
    \ifstrempty{#4}{}{\par\vspace{3mm}{\color{pfRamka}\hrule height 0.5pt}\vspace{2mm}%
      {\footnotesize #4}}
  \end{tcolorbox}}
\renewcommand\pfGlavaNastrojka{\renewcommand\pfVkladkaSsylka{glava:\pf@g}\colorlet{pfVkladkaCvet}{pfBrend}}
\renewcommand\pfGlavaVid[3]{\pf@plashka{#3}{#1}{#2}{}}
% Компактная плашка — когда первая тема главы не влезает под полную
\renewcommand\pfGlavaVidMal[3]{%
  \begin{tcolorbox}[enhanced, colback=white, colframe=pfRamka, boxrule=0.7pt, arc=2.2mm,
    left=4mm, right=4mm, top=1.6mm, bottom=1.6mm, before skip=0pt, after skip=2.5mm,
    underlay={\begin{scope}\clip[rounded corners=2.2mm] (frame.south west) rectangle (frame.north east);
      \draw[pfKletka, step=5mm, line width=0.4pt] (frame.south west) grid (frame.north east);\end{scope}}]
    {\color{pfVkladkaCvet}\pfBrendZhirnyj\fontsize{15}{15}\selectfont #1}\hspace{2.5mm}%
    {\pfBrendZhirnyj\fontsize{10.5}{12}\selectfont #2}\hfill
    {\color{pfSeryj}\pfBrendShrift\fontsize{7.5}{9}\selectfont #3}
  \end{tcolorbox}}
\renewcommand\pfZadanieVid[1]{\renewcommand\pfVkladkaSsylka{zadanie:#1}\colorlet{pfVkladkaCvet}{pfPurple}%
  \def\pf@sp{{\color{pfSeryj}\scriptsize Типы задания — по классификации банка задач Профиматики
    \hfill разбор\hspace{6mm}в банке}\par\vspace{1mm}}%
  \ifcsdef{pf@zadtipy@#1}{\edef\pf@tmp{\csuse{pf@zadtipy@#1}}%
    \expandafter\forcsvlist\expandafter\pf@dobavTip\expandafter{\pf@tmp}}{}%
  \pf@plashka{Практика · задание #1}{№\,#1}{\ifcsdef{pf@zadname@#1}{\csuse{pf@zadname@#1}}{Задание #1}}{\pf@sp}}
\newcommand\pf@dobavTip[1]{\appto\pf@sp{\pf@tipStroka{#1}}}
% строка типа в плашке задания: код · название · страница разбора · ссылка в банк
\newcommand\pf@tipStroka[1]{\edef\pf@y{\csuse{pf@tl@#1}}%
  \par\noindent\makebox[8mm][l]{\color{pfPurple}\pfMedium #1}%
  \ifcsdef{pf@have@tip:\pf@y}%
    {\parbox[t]{0.62\linewidth}{\raggedright\hyperlink{tip:\pf@y}{\csuse{pf@title@#1}}}\hfill
     \hyperlink{tip:\pf@y}{\makebox[10mm][r]{\color{pfAkcent}с.\,\getpagerefnumber{tip:\pf@y}}}}%
    {\parbox[t]{0.62\linewidth}{\raggedright\textcolor{pfSeryj}{\csuse{pf@title@#1}}}\hfill\makebox[10mm]{}}%
  \hspace{2mm}\href{\csuse{pf@bankurl@#1}}{\makebox[13mm][r]{\color{pfBrendTemn}задачи\,$\nearrow$}}%
  \par\vspace{0.4mm}}
% Заголовок типа в практике: полоса + ссылка на задачи этого типа в банке
\renewcommand\pfTipVid[1]{\par
  \begin{tcolorbox}[enhanced, colback=pfRazborFon, colframe=pfRazborFon, boxrule=0pt, arc=1.6mm,
    left=2.5mm, right=2.5mm, top=1.6mm, bottom=1.4mm, before skip=3mm, after skip=0mm]
    \pf@mark{tip}{\pf@y}%
    {\color{pfPurple}\pfBrendZhirnyj\fontsize{10}{12}\selectfont #1}\enspace
    {\pfBrendZhirnyj\fontsize{10}{12}\selectfont\csuse{pf@title@#1}}\hfill
    \href{\csuse{pf@bankurl@#1}}{\color{pfBrendTemn}\footnotesize задачи с разборами в банке\,$\nearrow$}
  \end{tcolorbox}}

% ======================================================================
% 6. Обложка, часть, оглавление
% ======================================================================
% Обложка — готовая картинка на всю страницу + кликабельные QR-коды
\newcommand\pfOblozhka{\thispagestyle{pfPustaya}%
  \begin{tikzpicture}[remember picture, overlay]
    \node[anchor=north west, inner sep=0pt] at (current page.north west)
      {\includegraphics[width=150mm]{стиль/бренд/обложка.png}};
    \begin{scope}[shift={(current page.north west)}]
      \node[anchor=north west, inner sep=0pt] at (13mm,-108.8mm)
        {\href{\pfVK}{\phantom{\rule{24mm}{23.7mm}}}};
      \node[anchor=north west, inner sep=0pt] at (13mm,-142.5mm)
        {\href{\pfYouTube}{\phantom{\rule{24mm}{23.7mm}}}};
      \node[anchor=north west, inner sep=0pt] at (13mm,-175.8mm)
        {\href{\pfTelegram}{\phantom{\rule{24mm}{23.5mm}}}};
    \end{scope}
  \end{tikzpicture}\null\clearpage}
\let\pfTitul\@gobbletwo   % отдельный титул не нужен — есть обложка

% Разделитель частей: \pfChast{1}{Теория}{пояснение}
\newcommand\pfChast[3]{\clearpage\thispagestyle{pfPustaya}%
  \renewcommand\pfNazvanieKnigi{Часть #1. #2}\renewcommand\pfKolontitul{Часть #1. #2}%
  \renewcommand\pfRazdel{}%
  \pdfbookmark[0]{Часть #1. #2}{bm-chast#1}%
  \begin{tikzpicture}[remember picture, overlay]
    \pfKletka{0}{0}{150}{212}
    \begin{scope}[shift={(current page.south west)}]
      \node[anchor=north east, inner sep=0pt] at (139mm,197mm) {\pfLogo{22mm}};
      \node[anchor=north west, inner sep=0pt, text width=118mm, align=left] at (11mm,150mm)
        {{\color{pfSeryj}\pfBrendShrift\fontsize{11}{13}\selectfont Часть #1}\par\vspace{2mm}
         {\setlength{\fboxsep}{1.6mm}\colorbox{pfBrend}{\color{white}\pfBrendZhirnyj\fontsize{30}{34}\selectfont #2}}\par\vspace{6mm}
         {\pfBrendShrift\fontsize{9}{13}\selectfont #3}};
    \end{scope}
  \end{tikzpicture}\null\clearpage}

% Оглавление: глава (крупно) и её темы (мелко), всё кликается
\renewcommand\pf@toc[3]{% ссылка, текст, метка страницы
  \par\noindent\hyperlink{#1}{#2}\dotfill\hyperlink{#1}{\color{pfSeryj}\pageref*{#3}}\par}
\newcommand\pf@tocMelko[2]{% вид, код
  \edef\pf@y{\csuse{pf@tl@#2}}%
  \ifcsdef{pf@have@#1:\pf@y}{%
  {\footnotesize\leftskip=11mm\rightskip=0pt plus 1fil\parfillskip=0pt\noindent
   \hyperlink{#1:\pf@y}{{\color{pfAkcent}\pfMedium #2}\enspace\csuse{pf@title@#2}}%
   \nobreak\leaders\hbox to 0.6em{\hss.\hss}\hskip 1em plus 1fill
   \hyperlink{#1:\pf@y}{\color{pfSeryj}\getpagerefnumber{#1:\pf@y}}\par}}{}}
\renewcommand\pfStrokaGlavy[1]{\edef\pf@g{\csuse{pf@gl@tl@#1}}%
  \vspace{2.2mm}\pf@toc{glava:\pf@g}{\makebox[11mm][l]{\pfChip{#1}}{\pfMedium\csuse{pf@gl@name@#1}}}{glava:\pf@g}%
  \vspace{0.8mm}\edef\pf@tmp{\csuse{pf@gl@temy@#1}}%
  \expandafter\forcsvlist\expandafter\pf@tocTema\expandafter{\pf@tmp}}
\renewcommand\pfStrokaZadaniya[1]{%
  \vspace{2.2mm}\pf@toc{zadanie:#1}{\makebox[11mm][l]{\pfChip[pfPurple]{№\,#1}}{\pfMedium
    \ifcsdef{pf@zadname@#1}{\csuse{pf@zadname@#1}}{Задание #1}}}{zadanie:#1}%
  \vspace{0.8mm}\edef\pf@tmp{\csuse{pf@zad@tipy@#1}}%
  \expandafter\forcsvlist\expandafter\pf@tocTip\expandafter{\pf@tmp}}
\newcommand\pf@tocTema[1]{\pf@tocMelko{tema}{#1}}
\newcommand\pf@tocTip[1]{\pf@tocMelko{tip}{#1}}
\newcommand\pfZagolovok[1]{{\pfBrendZhirnyj\fontsize{18}{22}\selectfont #1}\par\vspace{5mm}}
\newcommand\pfPodzagolovok[1]{\par\vspace{2mm}{\color{pfSeryj}\pfBrendShrift\fontsize{8}{10}\selectfont #1}\par\vspace{2.5mm}}

% Указатель «задание → темы»: карточка на задание, внутри — главы, справа номера тем
\newcommand\pf@ukChipSvet[1]{\tcbox[on line, colback=pfFon, colupper=pfAkcent, boxrule=0pt,
  arc=1mm, boxsep=0pt, left=1.2mm, right=1.2mm, top=0.75mm, bottom=0.6mm]%
  {\pfBrendZhirnyj\fontsize{8}{8}\selectfont #1}}
\renewcommand\pfUkazatel[2]{\par\noindent
  \begin{tcolorbox}[enhanced, colback=white, colframe=pfRamka, boxrule=0.6pt, arc=2mm,
    left=3mm, right=3mm, top=2mm, bottom=2mm, before skip=0pt, after skip=2.5mm,
    breakable=false]
    \noindent\pfChip[pfPurple]{№\,#1}\enspace{\pfMedium\csuse{pf@zadname@#1}}%
    \ifcsdef{r@zadanie:#1}{\hfill\hyperlink{zadanie:#1}{\footnotesize\color{pfPurple}разборы, с.\,\getpagerefnumber{zadanie:#1}}}{}%
    \par\vspace{1.2mm}{\color{pfRamka}\hrule height 0.4pt}\vspace{1mm}
    {\small #2}
  \end{tcolorbox}}
\renewcommand\pfUkGlava[2]{%
  \par\noindent\hangindent=10mm\hangafter=1\rightskip=0pt plus 1fil\parfillskip=0pt
  \makebox[10mm][l]{\pf@ukChipSvet{#1}}%
  \edef\pf@g{\csuse{pf@gl@tl@#1}}%
  \ifcsdef{r@glava:\pf@g}{\hyperlink{glava:\pf@g}{\csuse{pf@gl@name@#1}}}{\csuse{pf@gl@name@#1}}%
  \nobreak\leaders\hbox to 0.6em{\hss\textcolor{pfRamka}{.}\hss}\hskip 1em plus 1fill
  {\color{pfSeryj}\footnotesize темы}\ #2\par\vspace{0.4mm}}

% ======================================================================
% 7. Таблицы — единый стиль
% ======================================================================
% \begin{pfTablica}{lcc}            шапка (первая строка) + «зебра»
%   Длина & Площадь & Объём \\
%   ...
% \end{pfTablica}
% Опции: [bez shapki] — без шапки, только «зебра»; [levo] — прижать влево.
% Внутри: \pfPodshapka{текст} — строка-подзаголовок на всю ширину; \pfItog — итоговая строка.
% Ячейку с «&» внутри (aligned, cases) берите в фигурные скобки: {$...$} — иначе столбцы посчитаются неверно.
\definecolor{pfZebra}{RGB}{246,249,255}
\newif\ifpf@shapka \newif\ifpf@levo
\pgfkeys{/pfTab/.is family, /pfTab,
  bez shapki/.code={\pf@shapkafalse}, levo/.code={\pf@levotrue}}
\newcommand\pf@tabRamka{%
  \tikz\draw[pfRamka, line width=0.6pt, rounded corners=1.8mm] (1-|1) rectangle (last-|last);}
\NewDocumentEnvironment{pfTablica}{O{} m +b}{%
  \pf@shapkatrue\pf@levofalse\pgfkeys{/pfTab, #1}\pf@schitajStolbcy{#3}%
  \par\vspace{1.5mm}\begingroup\ifpf@levo\noindent\else\centering\fi\small
  \setlength{\tabcolsep}{2.2mm}%
  \ifpf@shapka
    \begin{NiceTabular}{#2}[rounded-corners=1.8mm, cell-space-limits=1.1mm,
      code-before={\rowcolor{pfFon}{1}\rowcolors{2}{white}{pfZebra}[respect-blocks]},
      code-after={\pf@tabRamka\tikz\draw[pfRamka, line width=0.6pt] (2-|1) -- (2-|last);}]
    \RowStyle[color=pfBrendTemn]{\pfMedium}#3
    \end{NiceTabular}%
  \else
    \begin{NiceTabular}{#2}[rounded-corners=1.8mm, cell-space-limits=1.1mm,
      code-before={\rowcolors{1}{white}{pfZebra}[respect-blocks]},
      code-after={\pf@tabRamka}]
    #3
    \end{NiceTabular}%
  \fi
  \par\endgroup\vspace{1.5mm}}{}
\newcommand\pfPodshapka[1]{\RowStyle[rowcolor=pfFon, color=pfBrendTemn]{\pfMedium\footnotesize}%
  \Block[l]{1-\pf@nCol}{#1}}
% число столбцов = максимум «&» в строке + 1 (считается по телу таблицы)
\ExplSyntaxOn
\seq_new:N \l__pf_rows_seq
\seq_new:N \l__pf_cells_seq
\int_new:N \l__pf_ncol_int
\cs_new_protected:Npn \pf@schitajStolbcy #1 {
  \int_set:Nn \l__pf_ncol_int {1}
  \seq_set_split:Nnn \l__pf_rows_seq {\\} {#1}
  \seq_map_inline:Nn \l__pf_rows_seq {
    \seq_set_split:Nnn \l__pf_cells_seq {&} {##1}
    \int_set:Nn \l__pf_ncol_int { \int_max:nn {\l__pf_ncol_int} {\seq_count:N \l__pf_cells_seq} } }
  \edef\pf@nCol{\int_use:N \l__pf_ncol_int}}
\ExplSyntaxOff
\newcommand\pfItog{\RowStyle{\pfMedium}}

% Схема Горнера: \pfGorner{корень}{коэффициенты через &}{получившиеся через &}{остаток}
\ExplSyntaxOn
\NewDocumentCommand\pfGorner{mmmm}{
  \seq_set_split:Nnn \l_tmpa_seq {&} {#2}
  \exp_args:Ne \pf@gornerTab { \seq_count:N \l_tmpa_seq } {#1}{#2}{#3}{#4}}
\ExplSyntaxOff
\newcommand\pf@gornerTab[5]{{\small\setlength{\arraycolsep}{2mm}%
  $\begin{NiceArray}{>{\color{pfBrend}}c|*{#1}{c}}[cell-space-limits=1mm,
    code-before={\rowcolor{pfFon}{2}}, rules/color=pfBrend!70, rules/width=0.7pt]
    & #3 \\ \hline
    #2 & #4 & \textcolor{pfBrend}{\mathbf{#5}}
  \end{NiceArray}$}}
